% RESEARCH_DRAFT_SHA256: 38517a78a3cda60b9c09cf2d227af35fd0ca2b5893ca3d60e5b4583ff01cfbe3
% Layout and source-reference wrappers expanded for the web renderer.
\documentclass[11pt]{article}
\usepackage[margin=25mm]{geometry}
\usepackage{amsthm}
\usepackage{fontspec,unicode-math}
\setmainfont{Latin Modern Roman}
\setmathfont{Latin Modern Math}
\usepackage{amsmath,amssymb,mathtools,xurl,hyperref}
\hypersetup{colorlinks=true,urlcolor=blue,linkcolor=blue}
\setcounter{secnumdepth}{0}
\setlength{\parindent}{0pt}
\setlength{\parskip}{5pt}
\setlength{\emergencystretch}{3em}
\newcommand{\sref}[2]{\hyperlink{source:#1:#2}{[S#2]}}
\newcommand{\cataloguescope}[1]{\paragraph{Source record.}#1}
\newcommand{\floor}[1]{\left\lfloor#1\right\rfloor}
\newtheorem{proposition}{Proposition}
\newtheorem{lemma}[proposition]{Lemma}
\begin{document}
\paragraph{Research manuscript.} OpenAI \texttt{gpt-6-astra}, reasoning \texttt{ultra}; 11 October 2026.
\paragraph{Preserved source body.} The complete catalogue mathematical text and sources follow with layout and source links expanded, including source qualifications and its historical preparation notes. New research begins at the section headed Research attempt.

\section{Robust Spectral Tensor Methods for Latent-Variable Learning}\label{statement}

\subsection{Definitions and mathematical statement}

\paragraph{Shared notation.}Standard number systems are $\mathbb N,\mathbb Z,\mathbb Q,\mathbb R,\mathbb C$. Specialized notation and index ranges follow the cited source and the definitions below.


A latent-variable model describes observations through hidden random variables. Tensor methods use low-order joint moments and a tensor factorization to infer the hidden parameters. Robustness concerns stability under noise or changes in those modeling assumptions.
\paragraph{Record scope and qualifications.}
The catalogue title is retained as metadata. Named certificates, contracts, bounds or frameworks in that title are not asserted as results of the original source extract.
The earlier extraction was flagged incomplete. The cleaned record gives a broad robustness direction without an error model, model family, or quantitative guarantee.

Source-record limitation: The supplied assessment flags an incomplete extraction; the cleaned qualitative robustness prompt still specifies no model or guarantee.
\paragraph{Statement recorded in the supplied source.}
Can spectral tensor-decomposition methods for learning latent Dirichlet allocation and related multi-view latent-variable models be made more robust to noisy data, statistical uncertainty, and incorrect modeling assumptions? \href{https://doi.org/10.4171/owr/2014/18}{[S2]}
\paragraph{Source status.}
The supplied catalogue labels this source record open; that metadata was checked against the supplied snapshot on 2026-10-11. The label is preserved for provenance and is not an independent verification of present open status or mathematical correctness. The catalogue literature assessment is dated 2026-08-21.
\subsection{Short English statement}
Improve the robustness of latent-variable tensor-learning methods.
\subsection{Sources}
\begin{itemize}
\item[S1] Ulam AI and the UnsolvedMath Contributors, UnsolvedMath: A Curated Collection of Open Mathematics Problems, catalogue entry. Catalogue curation and metadata: CC BY 4.0. The source extract is retained; definitions, English summaries and qualifications are editorial additions.
\url{https://huggingface.co/datasets/ulamai/UnsolvedMath}.
\item[S2] Mini-Workshop: Mathematical Physics meets Sparse Recovery (2015), Oberwolfach Reports, DOI 10.4171/owr/2014/18, pp. 1047--1073
\url{https://doi.org/10.4171/owr/2014/18}.
\item Source-specific Oberwolfach Reports licences, including CC BY-SA 4.0 where applicable, are retained; see THIRD\_PARTY\_NOTICES.md.
\end{itemize}
\paragraph{Preparation and review.}
The original source and UnsolvedMath contributors retain attribution for the source extract. Codex (GPT-6; exact serving revision not exposed) assisted with definitions, English summaries, scope notes and formatting. No new human mathematical authorship or independently verified solution is asserted. Catalogue identity checked against the complete supplied registry and current repository. Distinct upstream records, including repeated questions and broad or damaged extracts, are retained for complete coverage of the source CS category. A repeated record is not a new independent conjecture. Complete mathematical-equivalence and current-literature review has not been performed. Preparation and final review follow contribution prompts C and D, with source gaps disclosed. Structural and rendering checks do not establish a new theorem or open status.
\label{end:statement}

\section{Research attempt}
\subsection{Literature, source context and scope}
The full original catalogue body and source qualifications are preserved above. The relevant primary workshop passages were read in the cached publisher report. Historical claims and subsequent attributed results are distinguished from the mathematical work below. A restricted example, standard lemma or known theorem is not counted as new progress on the entire original target. Missing assumptions and repeated or rhetorical extracts are not silently turned into different solved problems.

\subsection{Whitening amplifies a nearly unidentifiable direction}
The source discusses multi-view latent-variable moments and asks for robustness to noise and misspecification. Consider a two-state latent variable with equal weights and three conditionally independent views whose conditional mean vectors are $u_1=(1,0)^T$ and $u_2=(1-\delta,\delta)^T$, where $0<\delta\le1$. These means can be realized by categorical views. The cross-view second moment is
\[
M=\tfrac12u_1u_1^T+\tfrac12u_2u_2^T,
\qquad \det M=\delta^2/4>0.
\]
Let $W=M^{-1/2}$ be symmetric whitening.
\begin{proposition}
Writing $\sigma$ for the smallest eigenvalue of $M$, one has $\sigma\le\delta^2/2$ and $\|W\|_2\ge\sqrt2/\delta$. An additive third-moment tensor error of Frobenius norm $\eta$ can become an error of norm $\eta\sigma^{-3/2}\ge2\sqrt2\eta/\delta^3$ after applying $W$ in every mode.
\end{proposition}
\begin{proof}
The Rayleigh quotient at $(0,1)^T$ is $\delta^2/2$, proving the eigenvalue bound. The eigenvalues of $W$ are the inverse square roots of those of $M$, giving its norm. Choose a unit eigenvector $v$ of $M$ for $\sigma$ and the perturbation $E=\eta v\otimes v\otimes v$. Its Frobenius norm is $\eta$, since the squared norm of a tensor product is the product of the squared factor norms. Applying $W$ in all modes gives $\eta\sigma^{-3/2}v\otimes v\otimes v$, proving the asserted amplification.
\end{proof}
Regularizing with $W_\lambda=(M+\lambda I)^{-1/2}$, $\lambda>0$, ensures $\|W_\lambda\|_2\le\lambda^{-1/2}$, but introduces a precise whitening defect. Diagonalizing $M$ shows
\[
\|I-W_\lambda M W_\lambda\|_2=\frac{\lambda}{\sigma+\lambda}.
\]
Thus reducing amplification by a fixed regularizer can produce a large bias as $\sigma$ tends to zero. Also, a second-moment perturbation $-\sigma vv^T$ of norm $\sigma$ makes the moment matrix singular, so ordinary inverse whitening is then undefined.
\subsection{Verification and remaining gaps}
The two eigenvalues, determinant and amplification were evaluated at $\delta=1/10$ and $1/100$, and the regularized defect was checked numerically. The argument concerns norm perturbations of moments; the displayed worst-direction third-moment error is not asserted to be the moments of another valid latent model. Accordingly it is not a statistical minimax impossibility theorem. The failed route is a conditioning-free stability bound for this whitening operation. Robust parameter recovery under a specified feasible noise or misspecification model remains unproved, and no general latent-variable-learning problem is settled.

\section{References and source attribution}
Primary workshop report: \url{https://doi.org/10.4171/owr/2014/18}. Relevant passages read: 2014/18 pp. 1061--1062, multi-view second/third moments, whitening and robustness questions.

All retrievals were checked on 11 October 2026. Original question attribution remains in the preserved Sources section. External theorems are attributed and not claimed as new discoveries or fully verified merely by citation.


\paragraph{Integration review.} Codex (GPT-6; exact serving revision not exposed) reviewed the written arguments and expanded layout and source-reference wrappers. Repository metadata and the production web renderer were checked using Python 3.13 and JavaScript with MathJax 3.2.2. This model-assisted review is not independent human proof verification; no PDF compilation or formal proof checking is asserted.

\section{Completion Estimate}
\noindent\textbf{COMPLETION: 0\%}
No substantiated new progress toward the entire original target is claimed. The worked results and failed route are delimited above; a known result, illustrative model, or auxiliary obstruction is not counted as original-target completion.

\section{Final status and attribution}
\textbf{Attempt status: unresolved.} This describes this research attempt, not a new ruling on the source catalogue's status or on external literature claims.
The mathematics and this self-audit were generated by OpenAI \texttt{gpt-6-astra}, reasoning \texttt{ultra}, on 11 October 2026. Codex (GPT-6; exact serving revision not exposed) orchestrated the batch, reviewed the written mathematical arguments, and contributed formatting/repository review. No human mathematical authorship, independent human proof checking, or formal verification is claimed. The historical preparation paragraphs in the reproduced catalogue body concern that body only.

Contribution prompt A was applied to the complete definition and supplied source record, followed by the adversarial checks of prompt B. Local reproduction files are \texttt{batch019/checks.py}, \texttt{computations.json}, and the per-ID audit JSON. They were not public links at preparation time. Numerical checks are finite and do not prove unrestricted results. No LaTeX engine was found locally; compilation was not performed by this research agent.

\end{document}
